\section{The sharp noiseless threshold for binary faces}\label{sec:noiseless-face-proof}
\begin{theorem}\label{thm:face-wHP}
In the noiseless binary model of Theorem~\ref{thm:box}, the probability
that all faces are resolved tends to one for fixed
\(A>C=3.4417465206\ldots\), and tends to zero for fixed \(0<A\le C\).
\end{theorem}

\begin{proof}
The recovery statement for \(A>C\) was already established in Theorem~\ref{thm:box}, so only the converse for \(0<A\le C\) remains. We prove it by applying the second moment method to the number of colliding pairs of opposite faces of a fixed codimension, showing that the expected number of collisions diverges while the variance is negligible compared to the squared mean.

Choose \(r\) to be the nearest even integer to \(\rho n\), where \(\rho\) is the maximizing fraction from Theorem~\ref{thm:box}. Each candidate pair consists of two opposite assignments on the same set of \(r\) prescribed coordinates, giving \(N=\binom nr2^{r-1}\) candidates in total. Orient each pair arbitrarily as a first and second assignment. For a cube word at distance \(d\) from the first assignment, its distance to the opposite assignment is \(r-d\); we define its signed jump in row \(j<r/2\) by \(g_j=\mathbf1_{d=j}-\mathbf1_{d=r-j}\). The two faces produce the same distance histogram, and hence collide, if and only if the sum of signed jumps over all sampled landmarks is zero in every row \(j<r/2\). Throughout the proof, write \(s=n^{1/3}\), \(W=n^{2/3}\), \(K=e^{As}\), \(v=\rho/4\), \(R=\sqrt{2vA}\), \(a=\tfrac23\sqrt{2v}A^{3/2}\), and \(\lambda_j=K\prob(\Bin(r,1/2)=j)\) for the expected landmark count in shell \(j\).

\step{Single candidate.}
We first estimate the common collision probability \(P_1\) of a single candidate and verify that the expected number of collisions \(NP_1\) tends to infinity even at \(A=C\). Because distinct shells are disjoint sets of cube words, \(P_1\) factors as a product over the opposing shell pairs \((j,r-j)\) for \(j<r/2\). Each opposing pair receives independent binomial landmark counts with mean \(\lambda_j\), and the probability that their difference is zero has negative logarithm \(\tfrac12(\log\lambda_j)_++O(1)\). Indeed, when \(\lambda_j\ge1\), the upper bound follows from the maximum atom bound for a binomial distribution and the lower bound follows from Cauchy--Schwarz on a Chebyshev interval around the mean; when \(\lambda_j<1\), requiring both shell counts to be zero yields the same estimate.

Restrict attention to the fixed window \(|j-r/2|\le MW\), choosing \(M\) large enough that the combined intensity of all shells outside this window is \(o(1)\). Applying the zero-count estimate to the shells outside the window contributes only \(o(1)\) to \(\log P_1\). Inside the window, Stirling's formula gives
\[
 \log\lambda_j=s\left(A-\frac{((j-r/2)/W)^2}{2v}\right)
                 -\tfrac12\log n+O(1).
\]
Rounding \(\rho n\) to the nearest even integer \(r\) changes the quadratic term by \(o(1)\) uniformly across the window. Summing the positive part of \(\log\lambda_j\) over the half band \(j<r/2\) yields
\[
 \log P_1=-an/2+(R/4)W\log n+O(W).
\]
To justify the error term, observe that the \(-\tfrac12\log n\) contribution shifts the boundary where \(\log\lambda_j\) becomes positive by \(O(s\log n)\) rows, which alters the sum approximation by at most \(O(s\log^2n)\). At the threshold \(A=C\) we have \(a=2h_{\rm face}(\rho)\), so the leading linear exponents in \(N\) and \(P_1\) cancel and the positive subleading term ensures that \(NP_1\) exceeds \(e^{c W\log n}\) for some \(c>0\). Below \(C\), the linear exponent is strictly positive, so \(NP_1\) grows exponentially in \(n\).

\step{Typical joint atoms.}
We next show that for almost all pairs of candidates, the joint collision probability is \((1+o(1))P_1^2\). For two candidates, let \(a_0,b_0\) count respectively the shared prescribed coordinates on which their first assignments agree and conflict, and let \(c\) count the prescribed coordinates exclusive to each candidate. Call the pair typical when \(c\ge\eta n\) and \(|a_0-b_0|/r\le n^{-1/12}\), for a sufficiently small fixed \(\eta>0\). Two independent uniform candidates fail to be typical with probability at most \(e^{-\Omega(n^{5/6})}\): the overlap size follows a hypergeometric distribution, which controls the event \(c\ge\eta n\), and conditional on the overlap, the signed agreement count obeys a Hoeffding bound when the oriented prescriptions are chosen independently and uniformly. Because neither condition depends on the orientations of the two pairs, the same probability bound holds under any arbitrary orientation.

To evaluate the joint collision probability of a typical pair, we pass to a Fourier integral, replace the Bernoulli weights by Poisson weights, and truncate to the retained window. Pairing each cube word with its bitwise complement negates both jump vectors simultaneously, which makes the exact joint characteristic function nonnegative: each complementary pair contributes a Bernoulli factor of the form \(1-2p(1-p)(1-\cos(\theta\cdot g))\) with \(p=K2^{-n}\). Uniformly over all frequencies, replacing these Bernoulli factors by the corresponding Poisson factors changes the integrand by a multiplicative factor \(\exp(O(2^np^2))=1+o(1)\), and the same comparison applies to the two marginal integrals. Next, projecting the jump vectors onto the retained window of rows changes the Poisson integrand by a uniform factor \(1+o(1)\), because the combined intensity of all discarded rows is \(o(1)\) and each summand in the log characteristic function has absolute value at most twice its intensity.

\step{Joint intensities and core localization.}
We now bound the cross-row intensities of a typical pair and localize the Fourier integrals to a small frequency core around the origin. For a typical pair, let \(I_{jk}\) denote the intensity of cube words that produce jumps in absolute folded rows \(j,k\) of the two blocks simultaneously. Uniformly over all retained rows,
\[
 I_{jk}\le \frac{\lambda_j\lambda_k}{K}\exp(o(s)).
\]
To prove this, observe that for a uniform word, the covariance matrix of its distances to the two first assignments has diagonal entries \(r/4\) and off-diagonal entry \((a_0-b_0)/4\). For each of the four sign choices corresponding to the folded rows, tilt each distance exponentially by its centered deviation divided by \(r/4\). Expanding the joint log moment generating function by Taylor's theorem over its independent coordinate blocks produces the sum of the two separate quadratic rate exponents plus a cross-covariance error bounded by \(O(n^{-1/12}s+1)\). The third-order remainder is \(O(n\,n^{-1})\), and comparing the resulting rate against the marginal probabilities via Stirling's formula adds \(O(\log n)\).

Call a retained row large if \(\lambda_j\ge e^{n^{1/4}}\). We claim that for each large row, there is a singleton jump affecting only that row in the first block (and no row in the second block) with intensity at least \(\lambda_j\exp[-O(n^{1/6}+\log n)]\). Indeed, conditional on the word's distance to the first assignment, freeze the three mismatch allocations on the shared agreeing, shared conflicting, and first-exclusive coordinates at their sequential hypergeometric modes; each step retains a polynomial fraction of the probability. Since the conditional slope of the second distance with respect to the first is \((a_0-b_0)/r\), the resulting conditional mean of the second query's distance differs from \(r/2\) by \(O(n^{7/12})\). The independent \(\Bin(c,1/2)\) mismatches on the second candidate's exclusive block then hit the neutral value \(r/2\) of the second query with the stated probability cost. This target value is an integer because \(r\) is even, and the required deviation on the exclusive block is \(o(c)\). Swapping the two blocks gives the analogous singleton construction for the second block.

These singleton jumps localize the Fourier integrals. Define the frequency core by the condition \(|\theta_j|\le e^{n^{1/5}}/\sqrt{\lambda_j}\) for every large row in both blocks. Outside the core, the joint integral is negligible relative to \(P_1^2\): if any large row violates the core bound, its singleton term alone bounds the integrand by \(\exp[-\exp(\Omega(n^{1/5}))]\), using the inequality \(1-\cos\theta\ge2\theta^2/\pi^2\) on the fundamental interval. The exact same localization holds for the two separate marginal integrals.

\step{Quadratic approximation and Gaussian comparison.}
Inside the core, we decouple the small rows from the large rows and replace the large-row exponents by their quadratic parts. The difference between the joint exponent and the sum of the two marginal exponents arising from a two-row jump that involves at least one small row is bounded by a constant times \(I_{jk}|\theta_j\theta_k|\). Summing over all such terms yields \(o(1)\) uniformly on the core: within each block, the sum of \(\lambda_j|\theta_j|\) over small rows is \(\le n e^{o(s)}\), whereas the sum over large rows is \(\le n\sqrt K e^{o(s)}\); applying the intensity bound for \(I_{jk}\) and dividing by \(K\) makes the total cross-contribution vanish. For the remaining large-row terms in both the marginal and joint exponents, we may replace the cosine terms by their quadratic approximations with a uniform \(o(1)\) error, because the sum of all quartic errors is at most \(O(n e^{4n^{1/5}-n^{1/4}})\). Here we bound the fourth power of a sum of two phases by eight times the sum of their fourth powers, and use that the total intensity touching row \(j\) is at most \(2\lambda_j\). As a consequence, the small-row integrals factor out identically in both the joint calculation and the product of marginals.

It remains to compare the large-row Gaussian integrals. For a given candidate, define the normalized radial Walsh vector of degree \(d\) as the sum of the degree-\(d\) coordinate characters on its prescribed support divided by \(\sqrt{\binom rd}\), weighted by the signs of the assignment. Vectors of distinct degrees are orthogonal. For two supports, the cross-Gram entry in degree \(d\) equals the signed sum over all \(d\)-subsets of the shared support, divided by \(\binom rd\). Its absolute value is at most \((1-\eta)^d\), and for each fixed \(d\) it equals \(((a_0-b_0)/r)^d+O_d(1/r)\) by comparing sampling with and without replacement. Because folded rows are antisymmetric under complementation, only odd degrees appear, so the entire cross-Gram matrix has Frobenius norm \(o(1)\) uniformly over typical pairs. Restricting to the large rows is a compression between orthonormal subspaces, which preserves this norm bound. Whitening each large-row coordinate by \(\sqrt{2\lambda_j}\) produces identity within-block covariances and a cross block \(T\) of Frobenius norm \(o(1)\), so the Gaussian determinant factor is \(\det(I-T^{\mathsf T}T)=1+o(1)\). In these whitened coordinates, all marginal variances of the Gaussian Fourier density are bounded, so the tails omitted outside the core are negligible. Combining these comparisons shows that every typical pair has joint collision probability \((1+o(1))P_1^2\) uniformly.

\step{Atypical contributions.}
Finally, we bound the contribution of atypical pairs by splitting them into three regimes according to how closely the two candidates align. First, suppose \(|a_0-b_0|/r\le1-\delta\) for a fixed small \(\delta\). In each block, choose one modal cell per folded row outside \(O_\delta(1)\) central rows; by the conditional mean contraction and the boundedness of rounding errors, any word in the modal cell of a row in one block has a strictly smaller folded radius in the other block. Consequently, the two collections of modal cells are disjoint, and ordering them by decreasing folded radius gives distinct leading coordinates, so the corresponding columns are independent. By the sequential-mode argument, each modal cell captures at least a polynomial fraction of its shell mass. Applying the conditional binomial atom bounds therefore yields a joint collision probability of at most \(P_1^2\exp(O_\delta(W\log n))\), which is overwhelmed by the atypical pair fraction \(\exp(-\Omega(n^{5/6}))\).

Second, suppose \(\sqrt n<r-|a_0-b_0|<\delta r\). The same modal-cell construction omits only \(O(\sqrt n)\) central rows, yielding the joint probability bound \(P_1^2e^{o(n)}\). Counting such pairs by the coordinate changes relative to a given assignment or its opposite shows that there are at most \(N e^{O(\delta\log(1/\delta))n}\) such configurations. For sufficiently small fixed \(\delta\), their total contribution is \(o(N^2P_1^2)\). Third, if \(r-|a_0-b_0|\le\sqrt n\), there are at most \(N e^{O(\sqrt n\log n)}\) such configurations, and the joint probability is trivially bounded by \(P_1\). Dividing their aggregate contribution by \(N^2P_1^2\) yields a quantity that vanishes even at \(A=C\), thanks to the positive subleading correction in \(NP_1\).

Combining the typical and atypical estimates shows that the variance of the number of opposite-face collisions is \(o((NP_1)^2)\). Chebyshev's inequality therefore establishes the converse.
\end{proof}
